% !TeX root = ../../Book1.tex
% 纲要标题：### 8.3 自由积
% 纲要位置：计划纲要.md，第 342 行。

\section{自由积}
\label{b1:sec:0803}

% 纲要内容（仅作写作计划，尚非教材正文）：
%
% * 自由积的约化字
% * 自由积的泛性质
% * 直积与自由积的区别
%

\subsection{自由积的约化字}
\label{b1:subsec:0803:01}

直积把两个群放在一起，并强制来自不同因子的元素交换。如果只想保留两个群各自已有的乘法，而不添加跨因子的关系，就应采用自由积。这与自由群相似，但一个字母现在可以是某个因子的任意非单位元。

\begin{definition}\label{b1:def:free-product-word}
设 \(G,H\) 为群，先把它们的非单位元作不相交标记。由这些元素组成的字称为自由积字；若相邻字母总来自不同因子，就称为约化字。空字也约化。
\end{definition}
从所有有限字出发，允许在同一因子内把相邻的 \(x,y\) 换为它们的乘积 \(xy\)，乘积为一时删去这两个字母；也允许这些操作的逆操作。由此生成的等价关系与串接相容。

\begin{lemma}\label{b1:lem:free-product-normal-form}
每个等价类中有唯一约化字。
\end{lemma}
\begin{proof}
从左到右读取字母，始终保存一个约化字。如果新字母与末尾来自不同因子，就接在末尾；若来自相同因子，就把末尾替换为两者乘积，乘积为一时删去末尾。这给出一个与原字等价的约化字。

验证输出不随允许的操作改变。设 \(x,y\) 来自同一因子 \(G\)，比较连续读取 \(x,y\) 与一次读取 \(xy\)。若原字为空或末尾来自 \(H\)，两次读取先添 \(x\) 再合并成 \(xy\)，与直接读取相同。若原末尾为 \(z\in G\)，且 \(zx\ne1\)，两次读取把末尾变为 \((zx)y=z(xy)\)；若 \(xy=1\)，这就恢复 \(z\)。若 \(zx=1\)，第一次删去 \(z\)，前一字母如存在必来自 \(H\)，第二次便接上 \(y\)；直接读取的结果由 \(z(xy)=y\) 给出，也一样。单位乘积的删除包含在这些情形中。交换 \(G,H\) 得到另一因子的同样结论。

因此一次合并或拆分同因子字母不改变算法输出。约化字的输出是自身，所以同一等价类不能包含两个不同约化字。
\end{proof}
等价类在串接下组成群：结合律从串接继承，空字为单位，逆由字母倒序取逆得到。称此群为 \(G,H\) 的\term[ziyou-ji-product]{自由积}[free product]，记为 \(G*H\)。长度为一的字分别给出 \(G,H\) 到 \(G*H\) 的单同态，因而可把原群看作其子群；两者的交只有单位元。
\symidx{\(G*H\)：两个群的自由积}

\ProseParagraphBreak
例如在 \(C_2*C_3\) 中，设 \(a^2=b^3=1\)，则
\[
(ab^2a)(ab)=ab^2b=a,
\]
但 \((ab)^k\) 对每个 \(k\ge1\) 都是长度 \(2k\) 的约化字，所以 \(ab\) 具有无限阶。自由积即使由有限群构成，也通常无限。

\subsection{自由积的泛性质}
\label{b1:subsec:0803:02}

\begin{theorem}\label{b1:thm:free-product-universal}
给定同态 \(f:G\to K\)、\(g:H\to K\)，存在唯一同态 \(u:G*H\to K\)，使它在两个因子上的限制分别为 \(f,g\)。
\end{theorem}
\begin{proof}
把一个字映为各字母像的乘积。同因子字母的合并不改变这个乘积，因为 \(f,g\) 各自保持乘法，所以映射在等价类上良定义，并保持串接乘法。两个因子共同生成 \(G*H\)，故这样的同态唯一。
\end{proof}
自由积的泛性质没有要求 \(f(G)\) 与 \(g(H)\) 交换。例如从 \(C_2*C_3\) 到 \(S_3\)，可以指定 \(a\mapsto(12)\)、\(b\mapsto(123)\)，得到满同态。若再令 \((ab)^2=1\)，商群就同构于 \(S_3\)：该关系给出 \(aba=b^{-1}\)，所以每个字化为 \(b^i a^\varepsilon\)，至多六个元素；相应置换像已经有六个元素。

\ProseParagraphBreak
与自由群相同，泛性质也确定了带指定因子映射的自由积，直到唯一同构。它还给出
\[
F(S\sqcup T)\cong F(S)*F(T)
\]
（这里取不相交副本）。具体地，任意指定 \(S\sqcup T\) 的像，等价于分别指定两个集合的像，再分别延拓为两个自由群上的同态。也可以在两个方向直接发送各个字母，利用生成性检查复合为恒等。

\ProseParagraphBreak
若 \(G=\langle S\mid R\rangle\)、\(H=\langle T\mid U\rangle\)，且生成符号互不相交，则
\[
G*H\cong\langle S\sqcup T\mid R\cup U\rangle.
\]
因为到任意群的同态，只须分别满足两个因子内部的关系。这个描述没有额外加入 \([s,t]=1\)；加入这些关系后才得到直积。

\subsection{直积与自由积的区别}
\label{b1:subsec:0803:03}

\begin{proposition}\label{b1:prop:product-versus-free-product}
自然满同态
\[
\pi:G*H\longrightarrow G\times H,\qquad
g\longmapsto(g,1),\quad h\longmapsto(1,h)
\]
的核是所有 \([g,h]\)（\(g\in G,h\in H\)）的正规闭包。
\end{proposition}
\begin{proof}
记该正规闭包为 \(N\)。跨因子交换子在直积中为单位，故 \(\pi\) 通过 \(Q=(G*H)/N\)。在 \(Q\) 中两个因子的像逐元素交换，因此
\[
G\times H\longrightarrow Q,\qquad
(g,h)\longmapsto \overline g\,\overline h
\]
是同态。它与 \(Q\to G\times H\) 的复合在两个因子上均为恒等，因子又生成相应群，所以两次复合都是恒等。由此 \(Q\cong G\times H\)，得到核的描述。
\end{proof}
当两个因子都非平凡时，选非单位元 \(g\in G,h\in H\)，字 \(ghg^{-1}h^{-1}\) 长四且约化，故自然同态不单射。直积中的相应交换子却必为单位。

\ProseParagraphBreak
特别地，\(C_2*C_2\) 是无限二面体群。设两个二阶生成元为 \(a,b\)，令 \(r=ab,s=a\)。则
\[
srs^{-1}=r^{-1},
\]
而 \(r\) 具有无限阶。每个交错字都能化为唯一的 \(r^k\) 或 \(r^ks\)，其中 \(k\in\mathbb Z\)。它可具体作用在整数直线上：令 \(a\) 作用为 \(m\mapsto-m\)，\(b\) 作用为 \(m\mapsto1-m\)，则 \(r\) 作用为平移 \(m\mapsto m-1\)。这同时展示了反射的几何来源与约化字的代数判别。对比 \(C_2\times C_2\)，后者只有四个元素，两次反射的乘积也只有二阶。
